\section{Limitações}\label{sec:limitacoes}

O agrupamento por descritores idênticos não cobre famílias de moléculas muito parecidas entre si, que continuam espalhadas entre treino e teste; as contagens de SMILES permitem decorar essas famílias. Dividir por família (esqueleto molecular) controlaria isso, mas exige uma biblioteca de química. Por isso a acurácia de \accGrupAmbos\% ainda é otimista para moléculas realmente novas.

O modelo escolhe sempre uma entre quatro classes: não existe a opção ``nenhuma das quatro'', então ele não sabe dizer que uma molécula não pertence a nenhum dos alvos. As 10 seeds variam o sorteio da divisão, mas os dados são os mesmos, de modo que o desvio padrão não mede a incerteza para dados novos, e não há validação externa. Por fim, as 238 moléculas sem descritores (0,86\%) foram descartadas, em sua maioria muito grandes, e as probabilidades da árvore não são calibradas.

Na comparação entre modelos, a busca da Random Forest e do \emph{gradient boosting} foi leve e feita só no cenário principal, de modo que o teto pode estar um pouco subestimado. A comparação entre as três árvores base e as curvas usam uma única seed. A Random Tree foi reproduzida no scikit-learn a partir da descrição do algoritmo do Weka, sem rodar o próprio Weka, e pode diferir dele em detalhes como o tratamento de empates. Os tempos de treino não são reportados porque foram medidos com a máquina ocupada por outros processos.

Na base de empréstimos, a principal limitação é a própria base. Distribuições quase uniformes e um rótulo que segue um degrau no score indicam dados sintéticos, gerados por uma regra; em dados de crédito reais, a relação entre os atributos e a aprovação tende a ser mais gradual, o que poderia favorecer a rede. A conclusão de que a árvore vence a rede vale, portanto, para esta base, e não para crédito em geral. A busca da rede cobriu 24 combinações e só o otimizador Adam, e a semente dos pesos iniciais mudou junto com a divisão dos dados, de modo que o desvio entre seeds mistura as duas fontes de variação. Os tempos foram medidos numa única máquina, num núcleo só; os valores absolutos mudam com o computador, e as razões entre modelos são mais estáveis que os tempos em si.
